\chapter{1980:Paul Berg and Walter Gilbert and Frederick Sanger}

\label{chap:chemistry1980}

The Nobel Prize in Chemistry for the year 1980 was awarded half to Paul Berg and half to Walter Gilbert and Frederick Sanger.

\textbf{Motivation:}

for his fundamental studies of the biochemistry of nucleic acids, with particular regard to recombinant-DNA (one half, Paul Berg) and for their respective contributions concerning the determination of base sequences in nucleic acids (the other half, jointly, Walter Gilbert and Frederick Sanger)

\section{Paul Berg}

\subsection*{Early Life and Education}

Paul Berg was born on 1926-06-30 in New York City, New York, USA.

\subsection*{Career and Major Research}

Berg studied at Pennsylvania State University and received his doctorate from Case Western Reserve University in 1952. He worked at the National Institutes of Health and at Stanford University, where he was professor of biochemistry. He developed methods for combining the DNA of different organisms and helped pioneer recombinant-DNA technology.

\subsection*{Nobel Prize-winning Work}

The work for which Paul Berg was awarded the Nobel Prize in Chemistry in 1980 was described by the Nobel Committee as: "for his fundamental studies of the biochemistry of nucleic acids, with particular regard to recombinant-DNA".

Berg constructed the first recombinant DNA molecules by splicing DNA from different species, creating the basis for genetic engineering.

\subsection*{Significance and Impact}

Recombinant DNA technology gave scientists the ability to isolate and amplify genes and proteins, transforming biology, medicine, and the biotechnology industry.

\subsection*{Later Career and Honors}

Berg remained at Stanford University until his retirement and later became a leading voice on the responsible use of genetic engineering. He died on 2023-02-15 in Stanford, California, USA.

\subsection*{Legacy}

Berg is regarded as one of the founders of genetic engineering and of the modern era of molecular biology.

\subsection*{Selected Publications}

"Biochemical Method for Inserting New Genetic Information into DNA of Simian Virus 40: Circular SV40 DNA Molecules Containing Lambda Phage Genes and the Galactose Operon of Escherichia coli" (Proceedings of the National Academy of Sciences, 1972)

\clearpage

\section{Walter Gilbert}

\subsection*{Early Life and Education}

Walter Gilbert was born on 1932-03-21 in Boston, Massachusetts, USA.

\subsection*{Career and Major Research}

Gilbert studied at Harvard University and received his doctorate from the University of Cambridge. He taught physics and later biology at Harvard University, where he was professor of molecular biology. With Allan Maxam he developed a chemical method for sequencing DNA.

\subsection*{Nobel Prize-winning Work}

The work for which Walter Gilbert was awarded the Nobel Prize in Chemistry in 1980 was described by the Nobel Committee as: "for their respective contributions concerning the determination of base sequences in nucleic acids".

Gilbert developed the Maxam-Gilbert method, in which DNA is labeled and cut with base-specific chemicals to determine its sequence, and he also discovered the mechanism of repression by which gene regulation is controlled.

\subsection*{Significance and Impact}

His DNA sequencing methods and his work on gene regulation were central to the rise of molecular biology and genomics.

\subsection*{Later Career and Honors}

Gilbert later worked in biotechnology and in the private sector before returning to Harvard. He is still alive.

\subsection*{Legacy}

Gilbert's sequencing methods helped make genome analysis possible and contributed to the early genomics industry.

\subsection*{Selected Publications}

"A New Method for Sequencing DNA" (Proceedings of the National Academy of Sciences, 1977)

\clearpage

\section{Frederick Sanger}

\subsection*{Early Life and Education}

Frederick Sanger was born on 1918-08-13 in Rendcomb, England.

\subsection*{Career and Major Research}

Sanger studied at the University of Cambridge, where he received his doctorate. He worked for his entire career at Cambridge, at the Department of Biochemistry and later at the Laboratory of Molecular Biology. He determined the amino acid sequence of insulin and later developed chain-termination methods for sequencing DNA.

\subsection*{Nobel Prize-winning Work}

The work for which Frederick Sanger was awarded the Nobel Prize in Chemistry in 1980 was described by the Nobel Committee as: "for their respective contributions concerning the determination of base sequences in nucleic acids".

Sanger developed the dideoxy chain-termination method, in which DNA synthesis is stopped at each base to read the sequence. He had earlier received the 1958 Nobel Prize for determining the structure of insulin.

\subsection*{Significance and Impact}

Sanger's sequencing method became the standard technology of genomics and was used to determine the human genome.

\subsection*{Later Career and Honors}

Sanger retired in 1983. He is the only person to have been awarded the Nobel Prize in Chemistry twice. He died on 2013-11-19 in Cambridge, England.

\subsection*{Legacy}

Sanger's methods for sequencing proteins and DNA laid the foundations of the genomics revolution.

\subsection*{Selected Publications}

"The Free Amino Groups of Insulin" (Biochemical Journal, 1945)

"DNA Sequencing with Chain-Terminating Inhibitors" (Proceedings of the National Academy of Sciences, 1977)

\clearpage

\section*{Chapter Summary}

The 1980 prize recognized Paul Berg, Walter Gilbert, and Frederick Sanger for their contributions to nucleic acid biochemistry. Their work established recombinant DNA technology and methods for determining DNA sequences.

